%!TEX root = ../main.tex
% =============================================================
%  CAPÍTULO 2 — METODOLOGÍA             (Entrega 2 · 5 % · 5/5)
% =============================================================

\chapter{Metodología y procedimientos}\label{cap:metodologia}

\begin{guia}[Guía de la cátedra: qué se evalúa en este capítulo]
La cátedra exige \textbf{dos metodologías}. El error más común es
entregar solo una.
\begin{enumerate}
    \item \textbf{Metodología de la investigación}: área temática,
          planteamiento y delimitación del problema, marco teórico
          previsto, diseño concreto (fases), indicadores medibles con
          fuente y umbral, técnicas de recolección, instrumentos, datos,
          procesamiento, análisis crítico y síntesis. Tipo de
          investigación: en general, aplicada.
    \item \textbf{Metodología de dirección del proyecto}: marco de trabajo
          justificado, alcance y entregables, cronograma con fechas reales
          (Gantt con hitos que coinciden con las entregas), costos, riesgos
          (referencia al capítulo~\ref{cap:riesgos}), calidad, recursos y
          satisfacción de las partes interesadas.
    \item \textbf{Declaración de uso de IA}: qué tareas fueron asistidas y
          cómo se verificaron. Los prompts se guardan en una carpeta
          compartida con la cátedra.
\end{enumerate}
Extensión orientativa: 6 a 10 páginas. Referencias habituales:
\citet{hernandezsampieri2014} para investigación y \citet{pmi2021}
para gestión.
\end{guia}

\section{Área temática y diseño de investigación}

El proyecto combinará desarrollo de software, análisis de datos
históricos y evaluación técnica reproducible. No será una investigación
exclusivamente empírica ni exclusivamente de ingeniería de software: la
validez de sus conclusiones dependerá tanto de que el sistema funcione
como de que sus estimaciones estén correctamente delimitadas y
evaluadas. El diseño será aplicado y de estudio de caso único, con AMBA
como área temática y CABA como unidad geográfica de validación.

\section{Planteamiento y delimitación}

La unidad principal de análisis del componente de aprendizaje automático
será un registro de siniestro con víctima en motocicleta. La salida del
modelo se interpretará, en todo momento, como gravedad condicionada a la
ocurrencia de ese siniestro, nunca como probabilidad de que un viaje
termine en accidente. Esta delimitación no será una aclaración retórica:
determinará qué preguntas podrá responder el sistema y cuáles quedarán
fuera de su alcance, incluida cualquier afirmación sobre seguridad de
una ruta o riesgo individual de una persona.

\section{Unidad de análisis, fuentes de datos e indicadores}

Las fuentes locales que alimentarán el pipeline de datos serán el
dataset de hechos de motocicleta ya depurado, complementado por el
dataset de siniestros viales con víctimas y cruzado con el dataset de
clima histórico. Este pipeline se aplica sobre los datos históricos que
entrenan y evalúan el modelo; cuando esos datasets históricos se
actualicen con nueva información oficial, se volverán a procesar con el
mismo pipeline para mantener la consistencia. Los aportes en vivo, como
noticias recientes y reportes de motociclistas, siguen un circuito de
ingesta y validación distinto, porque no pasan por el mismo proceso de
depuración estadística que los datos históricos.

El pipeline aplicará filtros de calidad: descarte de registros sin hora
o clima disponible, cruce por fecha y hora entre datasets, y tratamiento
de la edad de la víctima cuando cae fuera del rango 14-95 años
considerado válido. En ese último caso, en lugar de descartar el
registro o completar la edad con un valor inventado, se lo conserva
marcado con una categoría aparte (edad desconocida): la ausencia de ese
dato puede estar relacionada con la gravedad del siniestro, y ocultarla
sería un error metodológico mayor que dejarla explícita.

Sobre ese universo se derivarán hotspots y corredores viales. La capa de
inseguridad general se construirá sobre delitos de la Ciudad procesados
con clustering HDBSCAN (agrupamiento por densidad que no exige fijar de
antemano la cantidad de zonas), y se cruzará con los hechos de
motocicleta para estimar qué proporción de esa inseguridad general
coincide con la exposición del motociclista.

Las variables de entrada al modelo incluirán hora, día, comuna, tipo de
vía, contraparte, condición climática y campos demográficos de la
víctima cuando existan.

\section{Técnicas e instrumentos de recolección}

Se utilizarán datasets y servicios públicos documentados, scripts de
procesamiento en Python, pruebas automatizadas (pytest en el backend,
Vitest en el frontend) y control de versiones Git. El trabajo de campo
con motociclistas se realiza bajo protocolo, con consentimiento
informado y resultados verificables desde su diseño: se recogen sus
opiniones sobre el funcionamiento de la aplicación en rondas sucesivas,
bajo el mismo protocolo en cada una.

\section{Procesamiento, evaluación y criterios de validez}

Evaluar un modelo sobre un recorte al azar de los mismos años que usó
para entrenar suele mostrar un desempeño mejor al real, porque el modelo
puede terminar aprendiendo patrones propios de esa época en vez de
patrones que se sostengan en el tiempo. Por eso el proyecto va a dejar
por escrito, antes de reportar cualquier métrica, qué datos se usan para
entrenar, cuáles para ajustar decisiones intermedias sin tocar la
evaluación final, y cuáles quedan completamente reservados para medir el
desempeño una única vez. Ese documento, el contrato de evaluación, va a
organizar el tiempo en tres tramos: los años más antiguos del dataset se
usan para entrenar el modelo, un año completo posterior sirve para
comparar configuraciones y ajustar decisiones sin tocar la evaluación
final, y el año más reciente disponible queda reservado, sin tocarlo
hasta el final, para medir una única vez qué tan bien predice el modelo
sobre un período que nunca vio.

La imputación de edad, las categorías codificadas y el conjunto de
columnas se ajustarán exclusivamente con el conjunto de entrenamiento,
siguiendo la práctica documentada por scikit-learn (la biblioteca de
Python de referencia para aprendizaje automático clásico) para evitar
fuga de preprocesamiento entre partición de entrenamiento y partición de
evaluación \citep{sklearndocs}.

El desempeño del modelo se medirá con F1 macro (el promedio del
equilibrio entre precisión y recall calculado por separado para cada
clase de gravedad, sin favorecer a la clase mayoritaria) y con ECE
multiclase (error de calibración: qué tan cerca está la probabilidad que
informa el modelo de la frecuencia real con la que acierta); la
precisión y el recall por clase se derivarán de la matriz de confusión
correspondiente.

\section{Gestión del proyecto y cronograma}

La gestión del proyecto se centralizará en un repositorio de control de
versiones como fuente única de verdad: un backlog con identificadores
estables para hallazgos y oportunidades de mejora, una bitácora
cronológica de sesiones de trabajo, y un registro de decisiones que
distinguirá explícitamente una recomendación de una decisión
efectivamente tomada, evitando que una propuesta de mejora se confunda
con trabajo ya implementado.

El trabajo se acompañará de auditoría técnica periódica sobre el propio
código y la documentación, en lugar de concentrar esa revisión al final
del proyecto. Las decisiones tecnológicas relevantes (contenedores,
despliegue en la nube, proveedores externos) se tomarán pensando en la
operación sostenida del sistema y no solo en la demostración puntual de
una entrega.

\begin{figure}[H]
    \centering
    \begin{ganttchart}[
        hgrid, vgrid={*{6}{draw=none}, dotted},
        x unit=0.42mm,
        y unit chart=6mm, y unit title=7mm,
        time slot format=isodate,
        bar label font=\small, milestone label font=\small,
        title label font=\small,
        bar/.append style={fill=mainColor!60, draw=none},
        milestone/.append style={fill=mainColor, draw=none},
    ]{2026-03-16}{2026-11-20}
        \gantttitlecalendar{month=shortname} \\
        \ganttbar{Idea y resumen}{2026-03-16}{2026-04-10} \\
        \ganttbar{Metodología}{2026-04-06}{2026-05-05} \\
        \ganttbar{Introducción}{2026-05-06}{2026-06-09} \\
        \ganttbar{Marco teórico}{2026-06-01}{2026-07-28} \\
        \ganttbar{Marco tecnológico}{2026-07-01}{2026-08-25} \\
        \ganttbar{Justif. económica}{2026-08-10}{2026-09-15} \\
        \ganttbar{Riesgos}{2026-09-01}{2026-09-29} \\
        \ganttbar{Desarrollo del sistema}{2026-06-15}{2026-10-10} \\
        \ganttbar{Resultados}{2026-09-20}{2026-10-13} \\
        \ganttbar{Conclusiones}{2026-10-05}{2026-10-20} \\
        \ganttbar{Integración y revisión}{2026-10-21}{2026-11-16} \\
        \ganttmilestone{Entrega final}{2026-11-17}
    \end{ganttchart}
    \caption{Cronograma del proyecto (ciclo lectivo 2026)}
    \label{fig:gantt}
\end{figure}

\section{Declaración de uso de inteligencia artificial y trazabilidad}

Durante el proyecto se contó con el acompañamiento de herramientas de
inteligencia artificial en distintas tareas de apoyo: documentación
técnica, auditoría de código, y revisión y sugerencias de redacción
sobre partes de este informe, además de algunas intervenciones puntuales
de código bajo pedido y revisión previos a cada commit. Ningún dato,
cita o resultado incorporado a este documento se presenta sin
verificación en su fuente primaria o en la evidencia técnica del
repositorio; la IA no sustituye trabajo empírico no realizado.

\begin{table}[H]
    \centering
    \caption{Esquema de entregas y ponderación de la cátedra (ciclo lectivo 2026)}
    \label{tab:hitos}
    \small
    \begin{tabularx}{\textwidth}{l L l}
        \toprule
        \tb{Entrega} & \tb{Contenido} & \tb{Peso} \\
        \midrule
        1 & Idea del proyecto y abstract & 15\,\% \\
        2 & Capítulo 1. Introducción & 5\,\% \\
        3 & Capítulo 2. Metodología y procedimientos & 5\,\% \\
        4 & Capítulo 3. Marco teórico y tecnológico & 15\,\% \\
        5 & MVP & 5\,\% \\
        6 & Capítulo 4. Arquitectura y diseño del sistema, Presentación v1 & 15\,\% \\
        7 & Prototipo & 5\,\% \\
        8 & Capítulo 5. Justificación técnico-económica & 10\,\% \\
        9 & Capítulo 6. Análisis de riesgos, Presentación v2 & 5\,\% \\
        10 & Presentación final, proyecto funcional y defensa & 20\,\% \\
        \bottomrule
    \end{tabularx}
\end{table}
